\documentclass[10pt]{letter}

\usepackage[T1]{fontenc}
\usepackage[utf8]{inputenc}
\usepackage[margin=0.82in]{geometry}
\usepackage{hyperref}
\usepackage{enumitem}

\hypersetup{hidelinks}
\setlength{\parskip}{0.5em}
\setlength{\parindent}{0pt}

\begin{document}

\begin{flushright}
Yuhao Peng\\
School of Mathematics and Physics\\
China University of Geosciences (Wuhan)\\
Wuhan, Hubei 430074, China\\
\href{mailto:pngyuo@cug.edu.cn}{pngyuo@cug.edu.cn}
\end{flushright}

\vspace{1em}
8 November 2026

\vspace{1em}
Dear Editors,

\textbf{Re: Submission of ``Patient-adaptive conformal prediction sets for
reliable single-channel EEG sleep staging''}

We submit the enclosed manuscript for consideration as a research paper in
\emph{Biomedical Signal Processing and Control}.

Single-channel EEG sleep staging is increasingly accurate, but a point
prediction does not indicate which epochs need manual review. We reframe the
output as a set of admissible sleep stages and study whether the resulting
coverage guarantee remains useful when the unit of deployment is the patient
rather than the epoch. Using subject-wise folds on Sleep-EDF, we show that a
classifier with 0.898 pooled coverage leaves 45\% of patients below the
nominal 0.90 level and its worst patient at 0.774. This failure is invisible
to marginal coverage and to aggregate calibration error.

The manuscript makes four contributions:

\begin{enumerate}[leftmargin=1.2em,itemsep=0.08em,topsep=0.2em]
\item a patient-centred evaluation protocol for set-valued sleep staging,
including per-patient coverage, set size, review workload, and temporal
consistency;
\item a forward--backward posterior-deficit score that reduces the mean set
from 2.54 to 1.53 stages and the review workload from 76\% to 16\% at
comparable coverage;
\item per-patient adaptive conformal inference, with its regularization
chosen only on calibration subjects, which raises worst-patient coverage from
0.774 to 0.886 without using test labels for tuning;
\item an explicit negative result, with proofs, showing why dynamic-programming
sequence construction does not improve on the smoothed per-epoch score.
\end{enumerate}

The work is positioned at the interface of EEG signal processing and
trustworthy clinical decision support. It tests two neural backbones, one
20-subject cohort, an independent 78-subject cohort, a hardware-shift cohort,
and controlled signal corruptions. All results, the exact experimental
protocol, and executable code are publicly available at
\url{https://github.com/YuhaoPeng1103/SleepCP}.

The manuscript is original, is not under consideration elsewhere, and has not
been published previously. Both authors approved the submission. The authors
declare no competing interests and received no specific funding for this
work. We select the non-open-access publication route. We do not suggest
reviewers.

Thank you for considering the manuscript.

\vspace{1.5em}
Sincerely,

\vspace{1.5em}
Yuhao Peng\\
On behalf of Zhiwen Song and Yuhao Peng

\end{document}
